\scanpage{103}

\uline{Corollary} $N(u)=N(|u|)$\,; if $u=h+ik$ $N(h),N(k)\leq N(u)\leq N(h)+N(k)$, then $u\in L^N(E)\deleted{\unsure{F}}$ \deleted{$\Longleftrightarrow$}$\Longleftrightarrow h,k\in L^N(E)\deleted{\unsure{(F)}}$. \deleted{\unsure{and}}

\margin{Analogue for $L_N(E,F)$}

\uline{Proposition} $L^N(E,F)\subset L_0(E,F)$ if $N$ not equivalent to $N_\infty$\unsure{;} (otherwise $L^N(E,F)=L(E,F)$). Indeed, if $u\in L^N(E,F)$, \deleted{let us prove that}\inserted{then} $|u|\in L^N(E)$, let us prove $|u|$ compact (whence $u$ compact). Otherwise, there would exist a Hermitian operator $v\in L(E)$ such that $v|u|$ is an infinite Hermitian projection, but one would have $v|u|\in L^N(E)$. This is impossible, for one \deleted{\unsure{will have|would have} $(\rho_i)\in\ell^N$} would have $\sum|x_i|\deleted{\unsure{1}}\leq M$ for $N^\circ((x_i))\leq 1$, hence $N^\circ$ equivalent to $N_1$, which is absurd.

\uline{Corollary} $L^{N_p}(E,F)=L^p(E,F)$ for $1\leq p\leq+\infty$\,; \deleted{if $p<+\infty$, is}\inserted{if $p<+\infty$, is} even identical to $L_{N_p}(E,F)$. (for $\ell_N=\ell^N\unsure{=|\Longrightarrow}L_N(E,F)=L^N(E,F)$\unsure{)}

\uline{Proposition} $u\in L^N(E,F)$, $v\in L^{N^\circ}(F,G)\Rightarrow vu\in L^1(E,G)$

Hence Trace pairing between $L^N(E,F)$ and \deleted{\unsure{over $C$|$\to C$}} $L^{N^\circ}(F,E)$. By restriction, \deleted{$L^{N^\circ}(F,E)$ is a subspace of the dual} in principle, one has

\uline{Proposition} the dual of $L_N(E,F)$ is $L^N(E,F)$

\uline{Corollary 1} $L^N(E,F)$ and $L_N(E,F)$ are complete

\uline{Corollary 2} $L_N(E,F)$ reflexive $\Leftarrow$ \uline{$N$ reflexive}. \margin{\begin{tikzpicture}\draw[red,line width=0.8pt] (0,0)--(0.01,0.65);\end{tikzpicture}} Converse true if $E$ and $F$ of infinite dimension. \uline{In any case}, $L_N(E,F)$ reflexive $\Longleftrightarrow L^N(E,F)$ reflexive $\Longleftrightarrow\ldots$

\uline{Corollary 3} If $1<p<+\infty$, then $L^p(E,F)$ is a reflexive space having as dual $L^{p'}(F,E)$

\uline{Proposition} If $N,N',N''$ are Schatten norms such that $N(xy)\leq N'(x)N'(y)$, then $u\in L^{N'}(E,F)$, $v\in L^{N''}(F,G)$ implies $vu\in L^N(E,F)$, \inserted{and $N(vu)\leq N'(u)N''(v)$ \uline{Corollary}} If \uline{one of the \deleted{\unsure{either}} $N',N''$ is put at the} bottom, one can do likewise with the other. \margin{\begin{tikzpicture}\draw[red,line width=0.8pt] (0,0)--(0.04,1.15);\end{tikzpicture}}

(\deleted{by transposition, one reduces \unsure{to the second} \illegible{2}, by} \deleted{taking $N'''=N^\circ$,} By the second convexity theorem if $v$ and $u$ are \deleted{\unsure{there by contraction}, one} \deleted{to prove that if $N_1(xyz)\leq N'(x)N''(y)N'''(z)$, then} \unsure{compact.|compact.} One reduces to it easily, noting that \deleted{one has $N_1(wvu)\leq N'(u)N''(v)N'''(w)$.}

\[
N(u)=\lim N(\underline{pu}\deleted{\unsure{p|q}})=\lim_p N(pup)=\lim_p N(\unsure{up|u_p})
\]

$p$ orthogonal projection of finite rank, \uline{varying} in an increasing filtered set. (\unsure{For|Thus} $N(pvup)=N((pv)(up))\leq N''(pv)N^{\unsure{''|'}}(up)\leq N''(v)N^{\unsure{\circ|'}}(u)$). Finally, let us recall

\uline{Proposition} $N((\lambda_i u))\leq N(u)$

\uline{Proposition} Characterization of Schatten spaces.

\margin{104}
